\documentclass[10pt,a4paper]{article}
% Preamble of the manuscript on metric trial functions
\usepackage[T1]{fontenc}
\usepackage[utf8]{inputenc}
\usepackage{lmodern}
\usepackage[english]{babel}
\usepackage{microtype}
\usepackage[margin=24mm]{geometry}
\usepackage{amsmath,amssymb,amsthm,mathtools}
\usepackage{bm}
\usepackage{booktabs}
\usepackage{array}
\usepackage{xcolor}
\usepackage{caption}
\usepackage{tikz}
\usepackage{pgfplots}
\pgfplotsset{compat=1.18}
\usepgfplotslibrary{groupplots}
\usetikzlibrary{calc}
\usepackage[numbers,sort&compress]{natbib}
\usepackage[colorlinks=true,linkcolor=blue!50!black,citecolor=green!40!black,
            urlcolor=blue!50!black]{hyperref}
\usepackage[nameinlink,capitalise]{cleveref}

\captionsetup{font=small,labelfont=bf}
\linespread{1.08}

\newtheorem{proposition}{Proposition}
\newtheorem{lemma}[proposition]{Lemma}
\newtheorem{corollary}[proposition]{Corollary}
\theoremstyle{definition}
\newtheorem{definition}[proposition]{Definition}
\theoremstyle{remark}
\newtheorem{remark}[proposition]{Remark}
\crefname{proposition}{Proposition}{Propositions}
\crefname{lemma}{Lemma}{Lemmas}
\crefname{corollary}{Corollary}{Corollaries}
\crefname{definition}{Definition}{Definitions}
\crefname{remark}{Remark}{Remarks}

% notation
\newcommand{\PM}{\textsf{PM}}
\newcommand{\PP}{\textsf{PP}}
\newcommand{\PH}{\textsf{PH}}
\newcommand{\eE}{e_{E}}
\newcommand{\dd}{\,\mathrm{d}}
\newcommand{\bx}{\bm{x}}
\newcommand{\bxi}{\bm{\xi}}
\newcommand{\bs}{\bm{s}}
\newcommand{\R}{\mathbb{R}}
\newcommand{\Qp}{Q_p}
\newcommand{\Pp}{P_p}
\newcommand{\Tp}{T_p}
\newcommand{\Iref}{\mathcal{I}_{\mathrm{ref}}}
\newcommand{\Ilow}{\Pi_{\mathrm{low}}}
\newcommand{\T}{\mathsf{T}}
\newcommand{\data}{data}

% colours for p = 2..8 (dark = high p) and for the three trial spaces
\definecolor{p2}{HTML}{FDB863}
\definecolor{p3}{HTML}{F08B3E}
\definecolor{p4}{HTML}{D9541A}
\definecolor{p5}{HTML}{B2182B}
\definecolor{p6}{HTML}{7F1C6A}
\definecolor{p7}{HTML}{4A2A8A}
\definecolor{p8}{HTML}{1B2A5E}
\definecolor{cPP}{HTML}{555555}
\definecolor{cPM}{HTML}{C0392B}
\definecolor{cPH}{HTML}{1F5FA8}
\definecolor{cB1}{HTML}{C0392B}
\definecolor{cB2}{HTML}{1F5FA8}
\definecolor{cB3}{HTML}{2E8B57}
\definecolor{cB4}{HTML}{E08E0B}

\pgfplotsset{
  every axis/.append style={tick label style={font=\scriptsize}, label style={font=\small},
    title style={font=\small}, grid=major, grid style={gray!20}},
  conv/.style={
    xmode=log, ymode=log, x dir=reverse,
    xtick={1,0.5,0.25,0.125}, xticklabels={$1$,$\frac12$,$\frac14$,$\frac18$},
    xmin=0.1, xmax=1.25, unbounded coords=jump, mark size=1.2pt, line width=0.7pt,
  },
  pline/.style={color=p#1, mark=*, mark options={solid,fill=p#1}},
  sPP/.style={color=cPP, mark=o, dashed, mark options={solid}, line width=0.8pt, mark size=1.4pt},
  sPM/.style={color=cPM, mark=*, mark options={solid}, line width=0.8pt, mark size=1.4pt},
  sPH/.style={color=cPH, mark=square*, mark options={solid}, line width=0.8pt, mark size=1.3pt},
  meshplot/.style={width=3.3cm, height=3.3cm, axis equal image, hide axis,
    title style={font=\scriptsize, yshift=-1.2ex}},
}


\title{Well-conditioned bases and a hybrid metric--parametric trial space\\
       for unsymmetric spectral elements}
\author{Draft --- USEM-MATLAB, \texttt{studies/manuscript\_metric}}
\date{\today}

\begin{document}
\maketitle

\begin{abstract}
\noindent
Unsymmetric (Petrov--Galerkin) finite elements with parametric test and metric trial
functions reproduce polynomial displacement fields exactly on arbitrarily distorted meshes,
including meshes with $\det\bm J\le0$. For spectral elements of degree $p$ the metric trial
functions are obtained from a generalised Vandermonde system whose conditioning and whose
interpolation properties deteriorate rapidly with $p$ and with the distortion. We separate
the two effects and propose one remedy for each. (1)~\emph{Basis.} The trial space does not
depend on the polynomial basis used to set up the interpolation problem, but the condition
number $\kappa(\bm V)$ does. A tensor-product Lagrange basis on the Gauss--Lobatto--Legendre
points of the bounding box of the skew node coordinates makes $\bm V$ a permutation matrix on
affine elements ($\kappa=1$; metric and parametric functions coincide to round-off; the element
matrix is symmetric to round-off for every $p$), and on distorted elements reduces $\kappa(\bm V)$
to essentially the Lebesgue function of the interpolation at the reference nodes; compared with
monomials, $\kappa$ decreases by factors $4.6\cdot10^{3}$--$2.9\cdot10^{4}$ on distorted elements at
$p=16$, and the error of the computed shape functions by factors $3.6\cdot10^2$--$8.2\cdot10^3$. The Lebesgue constant itself is
basis invariant and cannot be improved this way. (2)~\emph{Hybrid space.} Completeness requires
only the complete polynomials $P_p(\bx)$ to be metric. The hybrid space $T_p=P_p(\bx)\oplus
\operatorname{span}\{L_a(\xi)L_b(\eta):a+b>p\}$ keeps metric (Cartesian) polynomials for
$P_p$ and writes the remaining functions in parametric coordinates. It is complete, frame
invariant, interpolatory, coincides with the conventional spectral element on bilinear
elements, and its interpolant is the conventional one plus a polynomial correction; we
derive a factorisation of its interpolation matrix that characterises unisolvence and a bound
of its Lebesgue constant in terms of the Gauss--Lobatto--Legendre value. In numerical
studies (single elements up to $p=16$, Jacobian-unlimited patch tests, and a thick cylinder on
meshes with fixed relative distortion, $p\le8$) the hybrid element keeps Lebesgue constants
near the parametric level where those of $Q_p$ in skew coordinates reach $10^{4}$--$10^{8}$,
reproduces polynomial fields to $\le10^{-9}$ also on folded meshes, has no eigenvalues with
negative real part and only small imaginary parts ($\le1.8\,\%$ of the modulus), and converges with the optimal order $p$ and quasi-optimally (discrete
error $0.78$--$1.00$ times the interpolation error) on all non-folded meshes, where the original
unsymmetric element loses optimality beyond moderate distortion and the conventional element
drops to order $\approx\lfloor p/2\rfloor$. Limitations (folded elements at high $p$, the
distortion dependence of the error constant) are documented.
\end{abstract}

\section{Introduction}\label{sec:intro}
Isoparametric elements lose accuracy on distorted meshes because the parametric shape
functions, composed with a non-affine element map, no longer span the complete polynomials
in Cartesian coordinates~\cite{Lee1993effects,Arnold2002Approximation}. The unsymmetric
finite element method of \citet{Rajendran2003novel} restores completeness by a
Petrov--Galerkin formulation: the test functions remain parametric, the trial functions are
\emph{metric}, i.e.\ polynomials in physical coordinates that interpolate the nodal values.
With the Jacobian determinant cancelled in the stiffness integral, the resulting elements
reproduce every displacement field of total degree $\le p$ exactly on arbitrarily distorted
meshes, also when $\det\bm J$ changes sign (Jacobian-unlimited elements, conditions C1--C3 of
\citet{Rajendran2010technique}). The approach has been developed for a wide range of element
types and applications, e.g.\
\cite{Ooi2004hexahedron,Ooi2007extension,Cen2012shape,Cen2015unsymmetric,Xie2016Modified,Shang2018airy,Huang2022Unsymmetric,Pfefferkorn2021mesh,Pfefferkorn2023hourglassing,Zaehringer2026On};
two known difficulties are the frame dependence of Cartesian metric functions, removed by
skew (corotational) coordinates~\cite{Ooi2008Remedies,Xie2016Modified}, and the
``interpolation failure'' of the metric functions on strongly distorted
elements~\cite{Ooi2008Remedies,Rajendran2010technique}.

For high-order (spectral) elements, the metric trial functions of degree $p$ in each
coordinate are obtained element by element from a generalised Vandermonde system with
$(p+1)^2$ unknowns. Two different properties of this construction degrade with $p$ and with
the element distortion, and they are often not distinguished:
\begin{enumerate}
\item the condition number $\kappa(\bm V)$ of the interpolation matrix, which controls the
  round-off in the computed shape functions and depends on the polynomial basis used to set up
  the system;
\item the Lebesgue constant $\Lambda$ of the interpolation in the trial space, which controls the
  approximation and the amplification of nodal perturbations; it is a property of the space and
  the node set and does not depend on the basis.
\end{enumerate}
A companion study of the thick cylinder of \citet[Sec.~5.9]{Rajendran2010technique} with
$p=2,\dots,8$ (\cref{sec:cylinder}) shows the practical consequence: on meshes whose relative
distortion does not decrease under refinement, the unsymmetric spectral element with metric
functions $Q_p$ in skew coordinates converges optimally only for moderate distortion; beyond, its
convergence is erratic, and the loss correlates with Lebesgue constants of $10^{3}$--$10^{8}$
rather than with the sign of $\det\bm J$.

This paper proposes one remedy for each property.
\begin{enumerate}
\item \emph{A well-conditioned basis} for the metric interpolation problem
  (\cref{sec:basis}): a tensor-product Lagrange basis on the Gauss--Lobatto--Legendre (GLL)
  points of the bounding box of the skew node coordinates. We prove that the trial space and its
  Lebesgue constant are basis invariant, that the new basis makes $\bm V$ a permutation matrix on
  affine elements, and that on distorted elements $\kappa_1(\bm V)$ equals the product of the
  reference Lebesgue function at the scaled nodes and the Lebesgue function of the metric
  interpolation at the reference points, i.e.\ the basis removes all ill-conditioning that is not
  inherent in the interpolation problem. It does \emph{not} reduce the Lebesgue constant.
\item \emph{A hybrid trial space} (\cref{sec:hybrid}) that uses metric polynomials only where
  completeness requires them, for $P_p(\bx)$, and parametric functions for the remaining
  $\tfrac12p(p+1)$ functions. We establish completeness, frame invariance and the reduction to
  the conventional spectral element on bilinear elements, derive a factorisation of the
  interpolation matrix that characterises unisolvence, and show that the hybrid interpolant is
  the GLL interpolant plus the interpolation error of a polynomial, which bounds its Lebesgue
  constant by the GLL value plus a term that vanishes on bilinear elements.
\end{enumerate}
\Cref{sec:numerics} supports the claims by element-level studies up to $p=16$, mesh-level
studies (smooth fields, eigenvalues), Jacobian-unlimited patch tests on the meshes of
\citet{Rajendran2010technique}, and the thick-cylinder study with $1512$ solutions.
\Cref{sec:conclusions} summarises the advantages and the limitations.

\section{Unsymmetric spectral elements}\label{sec:formulation}
\subsection{Petrov--Galerkin formulation}
Let $\hat\Omega=[-1,1]^2$ with the tensor grid of $(p+1)^2$ GLL points $\bxi_k$ and the
Lagrange polynomials $N_k\in Q_p(\bxi)$, $N_k(\bxi_l)=\delta_{kl}$. An element $\Omega^e$ is
the image of the isoparametric map $\bx(\bxi)=\sum_kN_k(\bxi)\bx_k$ with nodes $\bx_k$; $\bm
J=\partial\bx/\partial\bxi$. For linear elasticity with elasticity matrix $\bm D$ and strain
operator $\bm\varepsilon$, the unsymmetric element uses the parametric functions $N_k$ as test
functions and \emph{metric} functions $M_k$ as trial functions, which interpolate at the
element nodes, $M_k(\bx_l)=\delta_{kl}$, and span a space $S^e$. The element stiffness
matrix is
\begin{equation}
  \bm K^e=\int_{\hat\Omega}\bigl(\bm B_L^\ast\bigr)^{\T}\bm D\,\bm B_R\dd\xi\dd\eta ,
  \qquad \bm B_L^\ast=\bm B_L\det\bm J ,
  \label{eq:Ke}
\end{equation}
where $\bm B_L^\ast$ is built from $\operatorname{adj}\bm J\,\nabla_{\bxi}N_k$ and $\bm B_R$
from $\nabla_{\bx}M_k$ evaluated at $\bx(\bxi)$~\cite{Rajendran2010technique}. The load vector
is that of the conventional element. If $P_p(\bx)\subset S^e$ on every element (completeness),
the nodes are unisolvent for $S^e$, the metric functions are evaluated at the mapped quadrature
points, and loads and the (polynomial) integrand of~\eqref{eq:Ke} are integrated exactly, every
displacement field of total degree $\le p$ is reproduced exactly on any mesh, independently of the
sign of $\det\bm J$ (distortion immunity). The conventional spectral element (\PP) uses $S^e=\{N_k\circ
\bx^{-1}\}$ and $\bm B_L^{\T}\bm D\bm B_L\det\bm J$ as integrand.

The classical metric space is $Q_p$ in the corotational skew coordinates $\hat\bx=\bm
J_0^{-1}(\bx-\bx_0)$ ($\bm J_0$, $\bx_0$: Jacobian and image of the element centre), which are
affine in $\bx$ and frame invariant~\cite{Xie2016Modified}; hence $P_p(\bx)\subset Q_p(\hat
x,\hat y)$. We denote the resulting element by \PM. Given a basis $\{\phi_c\}_{c=1}^{n}$,
$n=(p+1)^2$, of $S^e$, the metric functions at query points $\bx_q$ follow from
\begin{equation}
  \bm V\bm M=\bm\Phi ,\qquad V_{ck}=\phi_c(\bx_k),\quad \Phi_{cq}=\phi_c(\bx_q),\quad
  M_{kq}=M_k(\bx_q),
  \label{eq:vandermonde}
\end{equation}
and analogously for the gradients.

\subsection{Two measures of the metric interpolation}
For nodal data $\bm f=(f_k)$ the metric interpolant is $\mathcal I_S\bm f=\sum_kf_kM_k$. Two
quantities characterise this construction.
\begin{itemize}
\item The \emph{condition number} $\kappa(\bm V)=\|\bm V\|\,\|\bm V^{-1}\|$ of the linear
  system~\eqref{eq:vandermonde}. With LU factorisation and partial pivoting the computed
  $\hat{\bm M}$ satisfies $\|\hat{\bm M}-\bm M\|/\|\bm M\|\lesssim c_n\,\kappa(\bm V)\,u$ with
  unit round-off $u\approx1.1\cdot10^{-16}$~\cite{Higham2002Accuracy}.
\item The \emph{Lebesgue constant}
  \begin{equation}
    \Lambda_S=\max_{\bx\in\Omega^e}\lambda_S(\bx),\qquad \lambda_S(\bx)=\sum_k|M_k(\bx)| ,
  \end{equation}
  the norm of $\mathcal I_S$ from $\ell^\infty$ nodal data to $L^\infty(\Omega^e)$. It bounds the
  interpolation error, $\|u-\mathcal I_Su\|_\infty\le(1+\Lambda_S)\inf_{v\in
  S^e}\|u-v\|_\infty$~\cite{Brutman1997Lebesgue}, and the amplification of perturbations of the
  nodal values (round-off, data errors) by the trial functions.
\end{itemize}
For the GLL tensor grid, $\Lambda_{\mathrm{GLL}}$ grows like $(\log p)^2$; it is $4.15$ at $p=8$
and $6.0$ at $p=16$. For the conventional element on any mesh $\Lambda_{\PP}=\Lambda_{\mathrm{GLL}}$,
since $\lambda(\bx(\bxi))=\sum_k|N_k(\bxi)|$.

\section{Improvement 1: a well-conditioned basis}\label{sec:basis}
\subsection{Basis invariance}
\begin{proposition}[basis invariance]\label{prop:invariance}
If $\bm V$ in~\eqref{eq:vandermonde} is non-singular, the metric functions are the unique
$M_k\in S^e$ with $M_k(\bx_l)=\delta_{kl}$; they, their Lebesgue function $\lambda_S$ and the
element matrix~\eqref{eq:Ke} do not depend on the basis $\{\phi_c\}$.
\end{proposition}
\begin{proof}
For another basis $\bm\psi=\bm C\bm\phi$ with non-singular $\bm C$, $\bm V_\psi=\bm C\bm
V_\phi$ and $\bm\Phi_\psi=\bm C\bm\Phi_\phi$, hence $\bm V_\psi^{-1}\bm\Phi_\psi=\bm
V_\phi^{-1}\bm\Phi_\phi$; the same holds for the gradients.
\end{proof}
Completeness, frame invariance, distortion immunity and unisolvence ($\bm V$ non-singular for
one basis iff for all) are therefore properties of the space $S^e$ and the node set. The basis
affects only the floating-point error, through $\kappa(\bm V)$. Conversely, \emph{no choice of
basis can reduce $\Lambda_S$}: approximation properties and the amplification of nodal
perturbations can only be improved by changing the space (\cref{sec:hybrid}).

\subsection{Candidate bases and scaling}
The original implementation uses monomials $\hat x^a\hat y^b$. For monomials on real nodes
symmetric about the origin, the condition number grows at least exponentially with the degree
\cite{Gautschi1988Lower}; on a parallelogram the 2D matrix is the Kronecker product of two 1D
Vandermonde matrices, so $\kappa_2(\bm V)$ is the square of the 1D value. We consider
tensor-product bases $\phi_a(\tilde s)\phi_b(\tilde t)$ with Legendre polynomials,
Chebyshev polynomials, or the Lagrange polynomials $\ell_a$ through the 1D GLL points. These
polynomials are bounded on $[-1,1]$ but grow rapidly outside, while the skew coordinates of
the nodes of a distorted element leave $[-1,1]^2$ (e.g.\ $|\hat x|\le1.27$ for the trapezoid of
\cref{sec:elements}). The node coordinates are therefore mapped affinely from their bounding
box onto $[-1,1]^2$,
\begin{equation}
  \tilde s=(\hat x-c_s)/r_s ,\qquad \tilde t=(\hat y-c_t)/r_t ,
  \label{eq:bbox}
\end{equation}
with centre $c$ and half-widths $r$ of the box. The map is affine and computed from frame
invariant quantities, so the space $Q_p(\hat x,\hat y)$ and frame invariance are preserved
(\cref{prop:invariance}).

\begin{lemma}[affine elements]\label{lem:affine}
On a parallelogram, the Lagrange basis on the GLL points in the scaled
coordinates~\eqref{eq:bbox} gives $\bm V=\bm\Pi$, a permutation matrix: $\kappa(\bm V)=1$,
$M_k=N_k\circ\bx^{-1}$ exactly, and $\bm K^e$ of~\eqref{eq:Ke} equals the symmetric matrix of the
conventional element.
\end{lemma}
\begin{proof}
For an affine map $\bx=\bx_0+\bm J_0\bxi$, $\hat\bx=\bxi$. The bounding box of the nodes is
$[-1,1]^2$, since the GLL grid contains the end points, so $(\tilde s,\tilde t)=\bxi$ and
$\phi_c(\bx_k)=\ell_a(\xi_k)\ell_b(\eta_k)=\delta_{ck}$ up to the numbering. Then $\bm M=\bm
\Pi^{\T}\bm\Phi$ is evaluated without arithmetic, $M_k(\bx(\bxi))=N_k(\bxi)$, and
$\nabla_{\bx}M_k=\bm J^{-\T}\nabla_{\bxi}N_k$, so $\bm B_L^\ast=\bm B_R\det\bm J$ and
$\bm K^e=\int\bm B_R^{\T}\bm D\bm B_R\det\bm J\dd\xi\dd\eta$.
\end{proof}

\begin{proposition}[distorted elements]\label{prop:kappa}
Let $\hat\bs_d$ be the reference GLL points in the scaled coordinates~\eqref{eq:bbox}, so that
$\phi_c(\hat\bs_d)=\delta_{cd}$, let $\tilde\bs_k$ be the scaled node coordinates, and let
$\lambda_{\mathrm{ref}}(\bs)=\sum_c|\phi_c(\bs)|$. Then
\begin{equation}
  (\bm V^{-1})_{kd}=M_k(\hat\bs_d),\qquad
  \kappa_1(\bm V)=\max_k\lambda_{\mathrm{ref}}(\tilde\bs_k)\;\max_d\lambda_S(\hat\bs_d).
  \label{eq:kappa1}
\end{equation}
\end{proposition}
\begin{proof}
With the query points $\bx_q=\hat\bs_d$, $\bm\Phi=\bm I$ and $\bm M=\bm V^{-1}$. The 1-norm is the
maximum absolute column sum: $\|\bm V\|_1=\max_k\sum_c|\phi_c(\tilde\bs_k)|$ and
$\|\bm V^{-1}\|_1=\max_d\sum_k|M_k(\hat\bs_d)|$.
\end{proof}
Since $\tilde\bs_k\in[-1,1]^2$, the first factor is bounded by the reference value
$\Lambda_{\mathrm{GLL}}$. With this basis, $\kappa(\bm V)$ is therefore the Lebesgue function of
the metric interpolation at the reference points of the bounding box, up to the factor
$\Lambda_{\mathrm{GLL}}$: the basis adds no ill-conditioning beyond that of the interpolation
problem itself. Note that the points $\hat\bs_d$ may lie outside $\Omega^e$, where the metric
polynomials extrapolate, so $\max_d\lambda_S(\hat\bs_d)$ can exceed $\Lambda_S$.

The Legendre and Chebyshev bases have no such identity and are not exact on affine elements
($\kappa_2=55$ and $2.6$ at $p=16$, respectively), but behave similarly on distorted elements
(\cref{sec:elements}). We therefore use the scaled Lagrange basis as the default of the metric
space $Q_p(\hat x,\hat y)$.

\section{Improvement 2: a hybrid metric--parametric trial space}\label{sec:hybrid}
\subsection{Rationale and definition}
Distortion immunity requires $P_p(\bx)\subset S^e$; nothing else. The space $Q_p(\hat x,\hat
y)$ satisfies this requirement with $\tfrac12p(p+1)$ additional functions of total degree $p+1$
to $2p$ in physical coordinates. On a non-affine element the nodes $\bx_k$, images of a tensor
grid, do not form a tensor grid in $(\hat x,\hat y)$, and interpolation with these high-degree
tensor terms on the perturbed grid drives $\Lambda_S$ up; \cref{prop:structure} below makes this
precise, and \cref{fig:lebesgue} shows that removing these terms brings $\Lambda$ back to the GLL
level. The
parametric functions, in contrast, interpolate on the exact GLL grid in $\bxi$. The hybrid
space keeps metric functions only for $P_p$ and uses parametric functions for the remainder.

\begin{definition}[hybrid trial space]\label{def:hybrid}
With the Legendre polynomials $L_a$,
\begin{equation}
  C_p=\operatorname{span}\{L_a(\xi)L_b(\eta):\ 0\le a,b\le p,\ a+b>p\},\qquad
  \tilde T_p^e=\{\,q\circ\bx+r:\ q\in P_p(\hat x,\hat y),\ r\in C_p\,\}.
  \label{eq:hybrid}
\end{equation}
The trial functions are the $M_k\in\tilde T_p^e$ with $M_k(\bxi_l)=\delta_{kl}$, evaluated
through $\bxi$: $u_h(\bx(\bxi))=\sum_ku_kM_k(\bxi)$.
\end{definition}
$C_p$ is the $L^2(\hat\Omega)$-orthogonal complement of $P_p(\bxi)$ in $Q_p(\bxi)$, so
$Q_p(\bxi)=P_p(\bxi)\oplus C_p$; this fixes the complement canonically. (Another complement
$C'$ of $P_p(\bxi)$ in $Q_p(\bxi)$, e.g.\ $\operatorname{span}\{\xi^a\eta^b:a+b>p\}$, gives the same
space only if $C'\subset P_p(\bx)\circ\bx\oplus C_p$, which holds when $P_p(\bxi)\subset\tilde
T_p^e$, e.g.\ on bilinear elements, but not in general.) In the
implementation $P_p$ is represented by $L_a(\tilde s)L_b(\tilde t)$, $a+b\le p$, in the scaled
skew coordinates~\eqref{eq:bbox}, and the interpolation matrix is
\begin{equation}
  \bm V_T=\begin{bmatrix}\bm Q\\ \bm R\end{bmatrix},\qquad
  Q_{jk}=q_j(\bx_k),\quad R_{ik}=r_i(\bxi_k),
  \label{eq:VT}
\end{equation}
with the $d=\tfrac12(p+1)(p+2)$ basis functions $q_j$ of $P_p$ and the $n-d$ functions $r_i$
of $C_p$. Gradients: $\nabla_{\bx}(q\circ\bx)=\nabla_{\bx}q$ exactly, $\nabla_{\bx}r=\bm
J^{-\T}\nabla_{\bxi}r$. We denote the element with~\eqref{eq:Ke} and~\eqref{eq:hybrid} by \PH.

\subsection{Properties}
\begin{proposition}\label{prop:hybrid}
\begin{enumerate}
\item[(a)] $P_p(\bx)\circ\bx\subset\tilde T_p^e$ for every element map: \PH\ is complete and
  distortion immune for fields of total degree $\le p$, independently of the sign of $\det\bm J$.
\item[(b)] $\tilde T_p^e$ is frame invariant and contains the rigid-body modes.
\item[(c)] If $\bx(\bxi)$ is bilinear (straight-sided element) and $\bm V_T$ is non-singular,
  then $\tilde T_p^e=Q_p(\bxi)$, $M_k=N_k$, and \PH\ is the conventional (symmetric) spectral
  element. In particular, on affine elements \PH, \PM\ and \PP\ coincide.
\item[(d)] The trial functions are interpolatory, so the GLL nodal-quadrature mass matrix is
  diagonal.
\end{enumerate}
\end{proposition}
\begin{proof}
(a)~by definition. (b)~$P_p(\hat x,\hat y)=P_p(\bx)$ is invariant under rigid motions, and so is
$\bxi$; $P_1(\bx)$ contains the rigid-body modes. (c)~For bilinear $\bx$, $x,y\in Q_1(\bxi)$, so
a monomial of total degree $j\le p$ in $(x,y)$ is a product of $j$ functions in $Q_1(\bxi)$ and lies
in $Q_j(\bxi)\subset Q_p(\bxi)$ \cite{Arnold2002Approximation}; $C_p\subset Q_p(\bxi)$. Hence
$\tilde T_p^e\subset Q_p(\bxi)$, with equality of dimensions if $\bm V_T$ is non-singular.
(d)~$M_k(\bxi_l)=\delta_{kl}$.
\end{proof}
By (c), \PH\ differs from the conventional element only on elements with curved edges (or
non-bilinear interior node placement); it is unsymmetric only there. A consequence of the same
argument is that the conventional Lagrange spectral element already contains $P_p(\bx)$ on
straight-sided quadrilaterals, so its loss of accuracy under distortion stems from curved
(non-bilinear) element maps, such as those of the mid-side node distortions studied in
\cref{sec:numerics}.

\paragraph{Unisolvence.} Let $\Iref$ denote interpolation in $Q_p(\bxi)$ at the GLL grid, and
$\Ilow$ the projection onto the Legendre coefficients with $a+b\le p$.
\begin{proposition}[factorisation]\label{prop:factor}
Let $\bm V_{\mathrm{ref}}$ be the Legendre Vandermonde matrix of the GLL grid (rows
$L_a(\xi_k)L_b(\eta_k)$, all $a,b\le p$) and $\bm A\in\R^{d\times d}$ the matrix of the map
$q\mapsto\Ilow\Iref(q\circ\bx)$ on $P_p$, $A_{ij}$ = coefficient of the $i$-th low-order Legendre
product in $\Iref(q_j\circ\bx)$. Then
\begin{equation}
  |\det\bm V_T|=|\det\bm V_{\mathrm{ref}}|\;|\det\bm A| ,
  \label{eq:factor}
\end{equation}
so the nodes are unisolvent for $\tilde T_p^e$ iff $\bm A$ is non-singular.
\end{proposition}
\begin{proof}
Order the rows of $\bm V_{\mathrm{ref}}$ as low-order ($a+b\le p$) rows $\bm L$ followed by the
rows $\bm R$ of~\eqref{eq:VT}. Each row of $\bm Q$ is a vector of nodal values and has a unique
expansion in the rows of $\bm V_{\mathrm{ref}}$: $\bm Q=\bm A^{\T}\bm L+\bm B\bm R$, the
coefficients being those of $\Iref(q_j\circ\bx)$. Hence
$\bm V_T=\begin{bmatrix}\bm A^{\T}&\bm B\\ \bm0&\bm I\end{bmatrix}\begin{bmatrix}\bm L\\ \bm R
\end{bmatrix}$.
\end{proof}
On an affine element $q\circ\bx\in P_p(\bxi)$, $\Iref$ is exact, $\Ilow$ is the identity on
$P_p(\bxi)$, and $\bm A$ is the matrix of the invertible composition with an affine map. Since
$\det\bm A$ depends continuously on the nodes, the nodes are unisolvent for every map
sufficiently close to an affine one. For strongly non-affine maps unisolvence is not
guaranteed; \cref{sec:cylinder} shows that it fails on some folded elements at $p\ge7$, exactly
where $\bm A$ becomes singular. Since $\bm V_{\mathrm{ref}}$ is fixed and well conditioned,
$\mathrm{rcond}(\bm A)$ is a cheap element-level indicator ($d\times d$ instead of $n\times n$).

\paragraph{Structure of the interpolant and Lebesgue constant.}
\begin{proposition}\label{prop:structure}
If $\bm A$ is non-singular, the hybrid interpolant of nodal data $\bm f$ is
\begin{equation}
  \mathcal I_T\bm f=\Iref\bm f+(\mathrm{Id}-\Iref)(q_{\bm f}\circ\bx),\qquad
  q_{\bm f}=\sum_jc_jq_j,\quad \bm A\bm c=\Ilow\Iref\bm f ,
  \label{eq:structure}
\end{equation}
i.e.\ $M_k=N_k+(\mathrm{Id}-\Iref)(q_k\circ\bx)$ with $q_k=q_{\bm e_k}$. Consequently
\begin{equation}
  \Lambda_T\le\Lambda_{\mathrm{GLL}}+\|\bm E_{\bx}\|\,\|\bm A^{-1}\|\,\gamma_p ,
  \qquad \bm E_{\bx}:\bm c\mapsto(\mathrm{Id}-\Iref)\Bigl(\sum_jc_jq_j\circ\bx\Bigr),
  \label{eq:lebesgue-bound}
\end{equation}
with $\gamma_p=\|\Ilow\Iref\|_{\ell^\infty\to\ell^\infty}$ depending on $p$ only and the norms
$\ell^\infty\to L^\infty(\hat\Omega)$ for $\bm E_{\bx}$. On bilinear elements $\bm
E_{\bx}=\bm0$ and $\Lambda_T=\Lambda_{\mathrm{GLL}}$.
\end{proposition}
\begin{proof}
Write $\mathcal I_T\bm f=q\circ\bx+r$. Interpolation at the GLL grid is exact on $Q_p(\bxi)\ni
r$, so $\Iref\bm f=\Iref(q\circ\bx)+r$, i.e.\ $r=\Iref\bm f-\Iref(q\circ\bx)$, which lies in $C_p$
iff $\Ilow\Iref(q\circ\bx)=\Ilow\Iref\bm f$, i.e.\ $\bm A\bm c=\Ilow\Iref\bm f$. Substituting $r$ gives
\eqref{eq:structure}; \eqref{eq:lebesgue-bound} follows from the triangle inequality and
$\|\Iref\|=\Lambda_{\mathrm{GLL}}$. For bilinear maps $q\circ\bx\in Q_p(\bxi)$
(\cref{prop:hybrid}(c)), so $(\mathrm{Id}-\Iref)(q\circ\bx)=0$.
\end{proof}
\Cref{prop:structure} states that the hybrid element is the conventional spectral element
plus a correction that restores completeness, and that this correction is an
\emph{interpolation error} of polynomials composed with the element map. It is small whenever
the map is close to bilinear: for $\bx=\bx_{\mathrm{bil}}+\delta\,\bm g$ with a smooth
perturbation $\bm g$, $\|\bm E_{\bx}\|=O(\delta)$, since $q\circ\bx_{\mathrm{bil}}\in
Q_p(\bxi)$. For curved elements with smooth edges $q\circ\bx$ is a polynomial of degree $p$
times the degree of the map, and $(\mathrm{Id}-\Iref)(q\circ\bx)$ decays with the smoothness of
the map. The same argument applies to $Q_p(\hat x,\hat y)$, with $q\in Q_p(\hat x,\hat y)$, $\Ilow$
replaced by the identity and an $n\times n$ matrix $\bm A_Q$; but then the correction contains the
interpolation errors of $\hat x^a\hat y^b$ up to total degree $2p$, which do not vanish even for
bilinear maps ($\hat x^p\hat y^p\circ\bx\in Q_{2p}(\bxi)\setminus Q_p(\bxi)$ in general). This is
consistent with \cref{fig:lebesgue}: on the straight-sided trapezoid $\Lambda_T=\Lambda_{\mathrm{GLL}}$,
whereas $\Lambda_Q$ grows to $4\cdot10^7$ at $p=16$.

\paragraph{Element matrix.} In~\eqref{eq:Ke} the columns of $\bm B_R$ that belong to $P_p$
are polynomial in $\bxi$, so the Jacobian cancellation is retained for them; the columns of the
complement contain $\bm J^{-1}$, and the corresponding integrand is
$\propto(\operatorname{adj}\bm J^{\T}\nabla_{\bxi}N)^{\T}\bm D(\operatorname{adj}\bm
J^{\T}\nabla_{\bxi}r)/\det\bm J$ (in strain-matrix form), rational as for curved isoparametric elements. Polynomial fields are
unaffected: for $\bm u\in P_p(\bx)$ the interpolant is exact, $\sum_ku_k\nabla M_k=\nabla\bm u$
is polynomial, and the integrand $(\bm B_L^\ast)^{\T}\bm D\bm\varepsilon(\bm u)$ is the same as for
\PM. On folded elements the complement part has a singular integrand where $\det\bm J=0$ inside
the element; \PH\ is undefined if $\det\bm J=0$ at a quadrature point. The cost per element
equals that of \PM\ plus one evaluation of $\bm J^{-1}$ per point; the parametric coordinates of
the quadrature points are known in all element routines.

\section{Numerical studies}\label{sec:numerics}
All computations use the USEM-MATLAB code base (GNU Octave~8.4, IEEE double precision), in
which the scaled Lagrange basis is the default of the metric space $Q_p(\hat x,\hat y)$
(\PM) and the hybrid space is available as a second trial space (\PH); \PP\ is the
conventional spectral element on the same nodes. GLL nodes are used throughout. The scripts
and data are listed in \cref{sec:reproduction}.

\subsection{Test elements and meshes}\label{sec:elements}
\emph{Single elements:} a parallelogram (affine), a trapezoid with straight edges (bilinear
map), an element with two curved edges (``curved'') and one with angular and curved-edge
distortion (``general''), $p=1,\dots,16$. \emph{Mesh level:} a $3\times2$ mesh on
$[0,3]\times[0,2]$ with curved-edge distortion $0.15$, $0.20$ and combined angular and
curved-edge distortion $0.30$ (minimum $\det\bm J=0.16$, $0.09$, $0.014$). \emph{Patch tests:}
the square panel of \citet[Secs.~5.5--5.7]{Rajendran2010technique}, $2\times2$ elements whose
internal mid-side nodes are rotated by $\theta=90^\circ$--$270^\circ$ (every element has
$\det\bm J\le0$ somewhere), with interior nodes from the bilinear map of the corners.
\emph{Convergence:} the thick cylinder of \citet[Sec.~5.9]{Rajendran2010technique} (quarter
annulus $1\le r\le3$, plane stress, $E=1000$, $\nu=0.25$, internal pressure $1000$),
$2N\times2N$ elements, $N=1,2,4,8$, generated with gmsh~\cite{Geuzaine2009Gmsh}; in each
$2\times2$ super-element the internal mid-side nodes are rotated by
$\alpha\in\{0^\circ,15^\circ,30^\circ,45^\circ,60^\circ,90^\circ\}$, so that the
\emph{relative} distortion is the same on all meshes and the element maps are not asymptotically
affine (\cref{fig:meshes}). Folded elements occur for $\alpha=45^\circ$ ($N\ge2$) and
$\alpha\ge60^\circ$. The arcs are represented exactly (Gordon--Hall blending~\cite{Gordon1973Transfinite}).
Three fields are computed: the Lam\'e solution, a manufactured polynomial of total degree~5, and
the analytic field $u_x=\tfrac1{10}(\sin3X\cos2Y+e^{XY})$, $u_y=\tfrac1{10}(\cos2X\sin3Y+Ye^{-X})$,
$X=x/3$, $Y=y/3$, with body forces and tractions derived from the exact fields.
Errors: relative energy norm $\eE$ (strains of the respective trial functions, signed
$\det\bm J$) and relative maximum nodal error $e_n$. Cases in which the code base finds the
interpolation matrix singular to working precision ($\mathrm{rcond}<\varepsilon$) are marked
``rej.''.

\begin{figure}[t]
  \centering
  \begin{tikzpicture}
    \begin{groupplot}[group style={group size=6 by 1, horizontal sep=1mm}, meshplot,
        xmin=-0.05, xmax=2.05, ymin=-1.05, ymax=1.05, width=3.0cm, height=3.0cm]
      \nextgroupplot[title={$\theta=90^\circ$}]
      \addplot[black, line width=0.3pt] table {\data/panel_mesh_N1_t090.dat};
      \nextgroupplot[title={$\theta=135^\circ$}]
      \addplot[black, line width=0.3pt] table {\data/panel_mesh_N1_t135.dat};
      \nextgroupplot[title={$\theta=175^\circ$}]
      \addplot[black, line width=0.3pt] table {\data/panel_mesh_N1_t175.dat};
      \nextgroupplot[title={$\theta=225^\circ$}]
      \addplot[black, line width=0.3pt] table {\data/panel_mesh_N1_t225.dat};
      \nextgroupplot[title={$\theta=270^\circ$}]
      \addplot[black, line width=0.3pt] table {\data/panel_mesh_N1_t270.dat};
      \nextgroupplot[title={$\theta=0^\circ$}]
      \addplot[black, line width=0.3pt] table {\data/panel_mesh_N1_t000.dat};
    \end{groupplot}
  \end{tikzpicture}\\[1mm]
  \begin{tikzpicture}
    \begin{groupplot}[group style={group size=6 by 1, horizontal sep=1mm}, meshplot,
        xmin=-0.05, xmax=3.05, ymin=-0.05, ymax=3.05, width=3.0cm, height=3.0cm]
      \nextgroupplot[title={$\alpha=0^\circ$}]
      \addplot[black, line width=0.3pt] table {\data/cyl_mesh_N2_a000.dat};
      \nextgroupplot[title={$\alpha=15^\circ$}]
      \addplot[black, line width=0.3pt] table {\data/cyl_mesh_N2_a015.dat};
      \nextgroupplot[title={$\alpha=30^\circ$}]
      \addplot[black, line width=0.3pt] table {\data/cyl_mesh_N2_a030.dat};
      \nextgroupplot[title={$\alpha=45^\circ$}]
      \addplot[black, line width=0.3pt] table {\data/cyl_mesh_N2_a045.dat};
      \nextgroupplot[title={$\alpha=60^\circ$}]
      \addplot[black, line width=0.3pt] table {\data/cyl_mesh_N2_a060.dat};
      \nextgroupplot[title={$\alpha=90^\circ$}]
      \addplot[black, line width=0.3pt] table {\data/cyl_mesh_N2_a090.dat};
    \end{groupplot}
  \end{tikzpicture}
  \caption{Top: panel meshes of the Jacobian-unlimited patch test (internal mid-side nodes
    rotated by $\theta$). Bottom: thick-cylinder meshes for $N=2$ ($4\times4$ elements); the finer
    meshes repeat the pattern on $N^2$ super-elements.}
  \label{fig:meshes}
\end{figure}

\subsection{Element level: conditioning and Lebesgue constants}\label{sec:num-elements}
\Cref{fig:basis} shows $\kappa_2(\bm V)$ and the error of the computed shape functions for
worst-case (unit) nodal data, $e_\delta=\max_{k,l}|M_k(\bx_l)-\delta_{kl}|$, for the space
$Q_p(\hat x,\hat y)$ in four bases. With monomials $\kappa_2$ grows by a factor $\approx5.7$
per degree already on the parallelogram ($3.5\cdot10^{11}$ at $p=16$), close to
$(1+\sqrt2)^2$, i.e.\ $(1+\sqrt2)^p$ for each 1D factor~\cite{Gautschi1988Lower}; the scaled
Lagrange basis gives $\kappa_2=1$ for every $p$ (\cref{lem:affine}), and the element matrix is
symmetric to round-off, whereas its asymmetry grows to $2\cdot10^{-6}$ with monomials
(\cref{tab:affine}). On the distorted elements the three alternatives lie within a factor~3 of
one another and are better than monomials by factors $4.6\cdot10^3$--$2.9\cdot10^4$ at $p=16$; $e_\delta$
follows $\kappa_2u$ and decreases by factors $3.6\cdot10^2$--$8.2\cdot10^3$ at $p=16$ (e.g.\
trapezoid: $3.0\cdot10^{-2}\to7.1\cdot10^{-6}$). For $p\le6$ all bases give $e_\delta<5\cdot10^{-13}$.
Without the bounding-box scaling~\eqref{eq:bbox} the orthogonal and Lagrange bases lose more than three
orders of magnitude on the trapezoid ($\kappa_2=2.2\cdot10^{15}$ instead of $8.6\cdot10^{11}$ at
$p=16$).

\begin{figure}[t]
  \centering
  \begin{tikzpicture}
    \begin{groupplot}[group style={group size=3 by 2, horizontal sep=13mm, vertical sep=16mm},
        width=5.2cm, height=4.4cm, ymode=log, xmin=1, xmax=16, xlabel={$p$},
        every axis plot/.append style={line width=0.7pt, mark size=1.1pt},
        unbounded coords=discard]
      \nextgroupplot[title={trapezoid}, ylabel={$\kappa_2(\bm V)$}, ymin=0.5, ymax=1e16,
        legend to name=leg:basis, legend columns=4, legend style={font=\footnotesize}]
      \addplot[cB1, mark=*] table[x=p, y=kappa_mon]{\data/metric-basis-trapezoid.dat};
      \addplot[cB2, mark=square*] table[x=p, y=kappa_leg]{\data/metric-basis-trapezoid.dat};
      \addplot[cB3, mark=triangle*] table[x=p, y=kappa_che]{\data/metric-basis-trapezoid.dat};
      \addplot[cB4, mark=diamond*] table[x=p, y=kappa_lag]{\data/metric-basis-trapezoid.dat};
      \legend{monomial, Legendre (scaled), Chebyshev (scaled), Lagrange (scaled)}
      \nextgroupplot[title={curved}, ymin=0.5, ymax=1e16]
      \addplot[cB1, mark=*] table[x=p, y=kappa_mon]{\data/metric-basis-curved.dat};
      \addplot[cB2, mark=square*] table[x=p, y=kappa_leg]{\data/metric-basis-curved.dat};
      \addplot[cB3, mark=triangle*] table[x=p, y=kappa_che]{\data/metric-basis-curved.dat};
      \addplot[cB4, mark=diamond*] table[x=p, y=kappa_lag]{\data/metric-basis-curved.dat};
      \nextgroupplot[title={general}, ymin=0.5, ymax=1e16]
      \addplot[cB1, mark=*] table[x=p, y=kappa_mon]{\data/metric-basis-general.dat};
      \addplot[cB2, mark=square*] table[x=p, y=kappa_leg]{\data/metric-basis-general.dat};
      \addplot[cB3, mark=triangle*] table[x=p, y=kappa_che]{\data/metric-basis-general.dat};
      \addplot[cB4, mark=diamond*] table[x=p, y=kappa_lag]{\data/metric-basis-general.dat};
      \nextgroupplot[ylabel={$e_\delta$}, ymin=1e-17, ymax=1]
      \addplot[cB1, mark=*] table[x=p, y=eKron_mon]{\data/metric-basis-trapezoid.dat};
      \addplot[cB2, mark=square*] table[x=p, y=eKron_leg]{\data/metric-basis-trapezoid.dat};
      \addplot[cB3, mark=triangle*] table[x=p, y=eKron_che]{\data/metric-basis-trapezoid.dat};
      \addplot[cB4, mark=diamond*] table[x=p, y=eKron_lag]{\data/metric-basis-trapezoid.dat};
      \nextgroupplot[ymin=1e-17, ymax=1]
      \addplot[cB1, mark=*] table[x=p, y=eKron_mon]{\data/metric-basis-curved.dat};
      \addplot[cB2, mark=square*] table[x=p, y=eKron_leg]{\data/metric-basis-curved.dat};
      \addplot[cB3, mark=triangle*] table[x=p, y=eKron_che]{\data/metric-basis-curved.dat};
      \addplot[cB4, mark=diamond*] table[x=p, y=eKron_lag]{\data/metric-basis-curved.dat};
      \nextgroupplot[ymin=1e-17, ymax=1]
      \addplot[cB1, mark=*] table[x=p, y=eKron_mon]{\data/metric-basis-general.dat};
      \addplot[cB2, mark=square*] table[x=p, y=eKron_leg]{\data/metric-basis-general.dat};
      \addplot[cB3, mark=triangle*] table[x=p, y=eKron_che]{\data/metric-basis-general.dat};
      \addplot[cB4, mark=diamond*] table[x=p, y=eKron_lag]{\data/metric-basis-general.dat};
    \end{groupplot}
  \end{tikzpicture}\\[1mm]
  \ref{leg:basis}
  \caption{Space $Q_p(\hat x,\hat y)$ on the distorted test elements: condition number of the
    interpolation matrix (top) and error of the Kronecker property of the computed shape
    functions (bottom) for four bases.}
  \label{fig:basis}
\end{figure}

\begin{table}[t]
  \centering\small
  \caption{Parallelogram (affine element), $Q_p(\hat x,\hat y)$: asymmetry
    $\|\bm K^e-\bm K^{e\T}\|_F/\|\bm K^e\|_F$ of the element matrix (zero in exact arithmetic) and
    condition number of $\bm V$.}
  \label{tab:affine}
  \begin{tabular}{@{}lcccccc@{}}
\toprule
$p$ & 4 & 6 & 8 & 10 & 12 & 16\\
\midrule
monomial & $2.2\cdot10^{-15}$ & $8.4\cdot10^{-14}$ & $4.2\cdot10^{-12}$ & $1.3\cdot10^{-10}$ & $2.9\cdot10^{-9}$ & $2.0\cdot10^{-6}$\\
Lagrange (scaled) & $1.6\cdot10^{-15}$ & $2.0\cdot10^{-15}$ & $4.6\cdot10^{-15}$ & $1.5\cdot10^{-14}$ & $1.5\cdot10^{-14}$ & $3.4\cdot10^{-14}$\\
\midrule
$\kappa_2(\bm V)$, monomial & $2.9\cdot10^{2}$ & $9.2\cdot10^{3}$ & $2.9\cdot10^{5}$ & $9.6\cdot10^{6}$ & $3.2\cdot10^{8}$ & $3.5\cdot10^{11}$\\
$\kappa_2(\bm V)$, Lagrange (scaled) & $1.0$ & $1.0$ & $1.0$ & $1.0$ & $1.0$ & $1.0$\\
\bottomrule
\end{tabular}

\end{table}

As \cref{prop:invariance} predicts, the reproduction of polynomials and the Lebesgue constant
are the same in all bases. \Cref{fig:lebesgue} compares $\Lambda$ of $Q_p(\hat x,\hat y)$ and of
the hybrid space $T_p$ on the three distorted elements. For $Q_p(\hat x,\hat y)$, $\Lambda$ grows
exponentially, to $4.0\cdot10^7$ (trapezoid), $5.5\cdot10^4$ (curved) and $6.7\cdot10^6$ (general) at
$p=16$. For $T_p$ it equals $\Lambda_{\mathrm{GLL}}$ on the trapezoid (bilinear map,
\cref{prop:structure}) and stays within $2\,\%$ (curved) and $40\,\%$ (general) of it ($6.1$ and
$8.5$ at $p=16$).

\begin{figure}[t]
  \centering
  \begin{tikzpicture}
    \begin{semilogyaxis}[width=8.5cm, height=5.6cm, xlabel={$p$}, ylabel={Lebesgue constant},
        xmin=2, xmax=16, ymax=1e9, legend columns=2,
        legend style={font=\scriptsize, at={(0.02,0.98)}, anchor=north west},
        every axis plot/.append style={line width=0.7pt, mark size=1.3pt}]
      \addplot[cPM, mark=*] table[x=p, y=LebQ_trapezoid]{\data/metric-basis-hybrid.dat};
      \addplot[cPH, mark=*, dashed] table[x=p, y=LebH_trapezoid]{\data/metric-basis-hybrid.dat};
      \addplot[cPM, mark=square*] table[x=p, y=LebQ_curved]{\data/metric-basis-hybrid.dat};
      \addplot[cPH, mark=square*, dashed] table[x=p, y=LebH_curved]{\data/metric-basis-hybrid.dat};
      \addplot[cPM, mark=triangle*] table[x=p, y=LebQ_general]{\data/metric-basis-hybrid.dat};
      \addplot[cPH, mark=triangle*, dashed] table[x=p, y=LebH_general]{\data/metric-basis-hybrid.dat};
      \addplot[black, dotted, line width=1pt] table[x=p, y=LebN]{\data/metric-basis-parallelogram.dat};
      \legend{$Q_p$ trapezoid, $T_p$ trapezoid, $Q_p$ curved, $T_p$ curved, $Q_p$ general,
        $T_p$ general, GLL}
    \end{semilogyaxis}
  \end{tikzpicture}
  \caption{Lebesgue constant of the metric interpolation in $Q_p(\hat x,\hat y)$ and in the hybrid
    space $T_p$ on the distorted test elements; GLL: parametric interpolation (any element).}
  \label{fig:lebesgue}
\end{figure}

\subsection{Mesh level: smooth field and eigenvalues}\label{sec:num-mesh}
On the distorted $3\times2$ meshes (\cref{fig:mesh}) the patch test of degree $p$ (harmonic
polynomial field) fails for the conventional element ($10^{-7}$--$10^{-2}$) and is passed by both
unsymmetric elements, with errors growing to $10^{-13}$--$10^{-12}$ at $p=10$ for $Q_p(\hat x,\hat
y)$ and $\le3.1\cdot10^{-15}$ for $T_p$. For the field $\bm u=(e^x\cos y,-e^x\sin y)$ all three elements
converge exponentially; for $p\le8$, $Q_p(\hat x,\hat y)$ is in most cases the more accurate of the
two unsymmetric elements (by factors up to $17$; exceptions: curved $0.20$, $p=6,8$; combined, $p=7$),
at $p=9,10$ $T_p$ is more accurate or within a factor $1.5$, and $T_p$ is $3.6$--$150$ times more accurate
than \PP\ for $p=5,\dots,10$ (curved $0.20$). The decisive difference
is in the spectrum of the generalised eigenvalue problem $\bm K\bm\phi=\lambda\bm M\bm\phi$ (consistent
mass, homogeneous Dirichlet conditions; reference: conventional element on the undistorted mesh,
$p=12$). With $Q_p(\hat x,\hat y)$, from $p=5$--$7$ on $2$--$25$ eigenvalues have negative real part,
the maximum of $|\operatorname{Im}\lambda|/|\lambda|$ approaches $1$, and the error of the eight
lowest eigenvalues stagnates at $10^{-5}$--$10^{-3}$ or reaches $O(1)$ (curved $0.20$, $p=6$: $1.9$).
With $T_p$ no eigenvalue has negative real part, $|\operatorname{Im}\lambda|/|\lambda|\le1.8\cdot10^{-2}$
over the whole spectrum, and the errors of the eight lowest eigenvalues do not exceed those of the
conventional element by more than $6\,\%$ and are smaller in most cases ($4.5\cdot10^{-7}$ for both at
$p=10$). An eigenvalue with negative real part
is an exponentially growing mode for every time integrator; for transient and modal analyses at
$p\gtrsim5$ on curved meshes this disqualifies $Q_p(\hat x,\hat y)$.

\begin{figure}[t]
  \centering
  \begin{tikzpicture}
    \begin{groupplot}[group style={group size=3 by 2, horizontal sep=15mm, vertical sep=17mm},
        width=5.2cm, height=4.4cm, ymode=log, xmin=1, xmax=10, xlabel={$p$},
        every axis plot/.append style={mark size=1.2pt}, unbounded coords=discard]
      \nextgroupplot[title={curved $0.20$}, ylabel={$e_n$, $e^x\cos y$},
        legend to name=leg:mesh, legend columns=3, legend style={font=\footnotesize}]
      \addplot[sPP] table[x=p, y=smooth_cla]{\data/metric-hybrid-curved20.dat};
      \addplot[sPM] table[x=p, y=smooth_skew]{\data/metric-hybrid-curved20.dat};
      \addplot[sPH] table[x=p, y=smooth_hyb]{\data/metric-hybrid-curved20.dat};
      \legend{parametric (PP), {$Q_p$, skew coordinates (PM)}, {$T_p$, hybrid (PH)}}
      \nextgroupplot[title={curved $0.20$}, ylabel={$e_\lambda$ (8 lowest)}]
      \addplot[sPP] table[x=p, y=eig_cla]{\data/metric-hybrid-curved20.dat};
      \addplot[sPM] table[x=p, y=eig_skew]{\data/metric-hybrid-curved20.dat};
      \addplot[sPH] table[x=p, y=eig_hyb]{\data/metric-hybrid-curved20.dat};
      \nextgroupplot[title={curved $0.20$}, ylabel={$\max|\operatorname{Im}\lambda|/|\lambda|$}, ymin=1e-4, ymax=2]
      \addplot[sPM] table[x=p, y=imag_skew]{\data/metric-hybrid-curved20.dat};
      \addplot[sPH] table[x=p, y=imag_hyb]{\data/metric-hybrid-curved20.dat};
      \nextgroupplot[title={angular + curved $0.30$}, ylabel={$e_n$, $e^x\cos y$}]
      \addplot[sPP] table[x=p, y=smooth_cla]{\data/metric-hybrid-combined30.dat};
      \addplot[sPM] table[x=p, y=smooth_skew]{\data/metric-hybrid-combined30.dat};
      \addplot[sPH] table[x=p, y=smooth_hyb]{\data/metric-hybrid-combined30.dat};
      \nextgroupplot[title={angular + curved $0.30$}, ylabel={$e_\lambda$ (8 lowest)}]
      \addplot[sPP] table[x=p, y=eig_cla]{\data/metric-hybrid-combined30.dat};
      \addplot[sPM] table[x=p, y=eig_skew]{\data/metric-hybrid-combined30.dat};
      \addplot[sPH] table[x=p, y=eig_hyb]{\data/metric-hybrid-combined30.dat};
      \nextgroupplot[title={angular + curved $0.30$}, ylabel={$\max|\operatorname{Im}\lambda|/|\lambda|$}, ymin=1e-4, ymax=2]
      \addplot[sPM] table[x=p, y=imag_skew]{\data/metric-hybrid-combined30.dat};
      \addplot[sPH] table[x=p, y=imag_hyb]{\data/metric-hybrid-combined30.dat};
    \end{groupplot}
  \end{tikzpicture}\\[1mm]
  \ref{leg:mesh}
  \caption{$3\times2$ meshes with curved-edge distortion: nodal error for the field
    $(e^x\cos y,-e^x\sin y)$ (left), error of the real parts of the eight eigenvalues of smallest
    modulus (centre), and largest relative imaginary part over the whole spectrum (right; zero for
    the symmetric conventional element).}
  \label{fig:mesh}
\end{figure}

\subsection{Jacobian-unlimited patch test}\label{sec:num-panel}
On the panel meshes the shear load case of \citet{Rajendran2010technique} has a cubic exact
solution, contained in both unsymmetric spaces for $p\ge3$. \Cref{tab:panel} gives the maximum
nodal error over the five distorted meshes (\cref{fig:meshes}). Both unsymmetric elements
reproduce the field on all meshes with folded and overlapping elements; the round-off level of
\PM\ reaches $1.1\cdot10^{-5}$ at $p=8$, that of \PH\ $6.7\cdot10^{-8}$. The conventional element has
errors of $13$--$140\,\%$. For $p=2$ (9-node element, corner-centroid interior node, $5\times5$ Gauss
points) \PH\ passes the constant and linear strain patch tests on all meshes and converges with
the energy-norm rates $2.02$ and $2.21$ on the folded mesh sequences B and C of the paper (PM-QUAD9$^\ast$:
$2.00$, $1.98$; conventional element: $0.70$, $2.50$, the latter from a negative displacement on
the coarsest mesh).

\begin{table}[t]
  \centering\small
  \caption{Jacobian-unlimited patch test (cubic field), maximum relative nodal error over the
    panel meshes with $\theta=90^\circ,\dots,270^\circ$.}
  \label{tab:panel}
  \begin{tabular}{@{}lcccccc@{}}
\toprule
 & $p=3$ & $p=4$ & $p=5$ & $p=6$ & $p=7$ & $p=8$\\
\midrule
parametric (\PP) & $1.6\cdot10^{-1}$ & $4.9\cdot10^{-1}$ & $1.5\cdot10^{-1}$ & $2.2\cdot10^{-1}$ & $1.3\cdot10^{-1}$ & $1.4$\\
$Q_p(\hat x,\hat y)$ (\PM) & $1.3\cdot10^{-13}$ & $3.3\cdot10^{-11}$ & $9.8\cdot10^{-11}$ & $6.5\cdot10^{-11}$ & $1.1\cdot10^{-9}$ & $1.1\cdot10^{-5}$\\
$T_p$ (\PH) & $6.4\cdot10^{-13}$ & $1.2\cdot10^{-11}$ & $1.6\cdot10^{-11}$ & $1.5\cdot10^{-10}$ & $2.7\cdot10^{-10}$ & $6.7\cdot10^{-8}$\\
\bottomrule
\end{tabular}

\end{table}

\subsection{Thick cylinder: accuracy, convergence and stability}\label{sec:cylinder}
\paragraph{Polynomial field.} \Cref{tab:cylpoly} gives the maximum nodal error for the degree-5
field over all meshes and $p=5,\dots,8$. Both unsymmetric elements are exact in exact arithmetic;
in floating point the error of \PM\ grows with the distortion to $1.1\cdot10^{-4}$ ($45^\circ$) and
$1.6\cdot10^{-2}$ (folded mesh $90^\circ$), that of \PH\ stays $\le9.9\cdot10^{-10}$ on every mesh on which
it can be built. With exact integration (\PM: $n_q\ge p^2+p/2$ points) the residual of the exact
nodal values is $9.9\cdot10^{-7}$ for \PM\ and $5.0\cdot10^{-11}$ for \PH\ ($\alpha=60^\circ$, $N=2$,
$p=8$), with $\sigma_{\min}/\sigma_{\max}(\bm K)=1.1\cdot10^{-15}$, $4.8\cdot10^{-8}$ and $2.9\cdot10^{-6}$ for
\PM, \PH\ and \PP.

\begin{table}[t]
  \centering\small
  \caption{Thick cylinder, polynomial field (degree 5): maximum relative nodal error over $N=1,\dots,8$
    and $p=5,\dots,8$; in parentheses the number of rejected cases (singular interpolation matrix).}
  \label{tab:cylpoly}
  \begin{tabular}{@{}lrrrrrr@{}}
\toprule
 & $\alpha=0^\circ$ & $\alpha=15^\circ$ & $\alpha=30^\circ$ & $\alpha=45^\circ$ & $\alpha=60^\circ$ & $\alpha=90^\circ$\\
\midrule
\PM & $1.6\cdot10^{-11}$ & $3.8\cdot10^{-10}$ & $7.6\cdot10^{-9}$ & $1.1\cdot10^{-4}$ & $1.2\cdot10^{-5}$ (4) & $1.6\cdot10^{-2}$ (3)\\
\PH & $4.3\cdot10^{-13}$ & $3.3\cdot10^{-13}$ & $5.1\cdot10^{-13}$ & $9.9\cdot10^{-10}$ & $2.8\cdot10^{-10}$ (3) & $9.7\cdot10^{-10}$ (3)\\
\bottomrule
\end{tabular}

\end{table}

\paragraph{Convergence.} \Cref{fig:cyl-h} shows the energy error of the analytic field and of
the Lam\'e field at $\alpha=30^\circ$, \cref{tab:rates} the rates on the finest pair of meshes. On
these non-asymptotically-affine meshes the conventional element converges with
$\approx\lfloor p/2\rfloor$, because the parametric space contains $P_k(\bx)$ only for $2k\le p$
\cite{Arnold2002Approximation}. \PM\ is optimal for $\alpha=15^\circ$, $p\le6$, but at
$\alpha=30^\circ$ only for $p\le3$; beyond, its error is erratic and can increase under refinement
(analytic field, $p=4$: $\eE=2.2\cdot10^{-4}$ at $N=4$, $4.3\cdot10^{-4}$ at $N=8$). \PH\ converges with the
optimal order $p$ for all $p$ at $15^\circ$ and $30^\circ$ (Lam\'e: $1.96$--$7.87$). At $N=8$,
$\alpha=30^\circ$ (analytic field), its energy error is $8$--$2.6\cdot10^4$ times smaller than that of \PP,
and $1.4$--$6.7\cdot10^4$ times smaller than that of \PM\ for $p\ge3$ (equal at $p=2$). Under
$p$-refinement (\cref{fig:cyl-p}) \PH\ converges exponentially to $3.5\cdot10^{-13}$ ($30^\circ$) and
$1.1\cdot10^{-11}$ ($45^\circ$, folded elements), \PM\ only for $\alpha\le15^\circ$.

\begin{figure}[t]
  \centering
  \begin{tikzpicture}
    \begin{groupplot}[group style={group size=3 by 2, horizontal sep=3mm, vertical sep=9mm,
        yticklabels at=edge left, xticklabels at=edge bottom}, conv, width=5.0cm, height=4.6cm,
        ymin=1e-14, ymax=1, xlabel={$h\propto1/N$}]
      \nextgroupplot[title={\PM, analytic}, ylabel={$\eE$},
        legend to name=leg:p, legend columns=7, legend style={font=\footnotesize}]
      \addplot[pline=2] table[x=h, y=eE2]{\data/cyl_h_smooth_PM_a030.dat};
      \addplot[pline=3] table[x=h, y=eE3]{\data/cyl_h_smooth_PM_a030.dat};
      \addplot[pline=4] table[x=h, y=eE4]{\data/cyl_h_smooth_PM_a030.dat};
      \addplot[pline=5] table[x=h, y=eE5]{\data/cyl_h_smooth_PM_a030.dat};
      \addplot[pline=6] table[x=h, y=eE6]{\data/cyl_h_smooth_PM_a030.dat};
      \addplot[pline=7] table[x=h, y=eE7]{\data/cyl_h_smooth_PM_a030.dat};
      \addplot[pline=8] table[x=h, y=eE8]{\data/cyl_h_smooth_PM_a030.dat};
      \legend{$p=2$, $p=3$, $p=4$, $p=5$, $p=6$, $p=7$, $p=8$}
      \nextgroupplot[title={\PH, analytic}]
      \addplot[pline=2] table[x=h, y=eE2]{\data/cyl_h_smooth_PH_a030.dat};
      \addplot[pline=3] table[x=h, y=eE3]{\data/cyl_h_smooth_PH_a030.dat};
      \addplot[pline=4] table[x=h, y=eE4]{\data/cyl_h_smooth_PH_a030.dat};
      \addplot[pline=5] table[x=h, y=eE5]{\data/cyl_h_smooth_PH_a030.dat};
      \addplot[pline=6] table[x=h, y=eE6]{\data/cyl_h_smooth_PH_a030.dat};
      \addplot[pline=7] table[x=h, y=eE7]{\data/cyl_h_smooth_PH_a030.dat};
      \addplot[pline=8] table[x=h, y=eE8]{\data/cyl_h_smooth_PH_a030.dat};
      \nextgroupplot[title={\PP, analytic}]
      \addplot[pline=2] table[x=h, y=eE2]{\data/cyl_h_smooth_PP_a030.dat};
      \addplot[pline=3] table[x=h, y=eE3]{\data/cyl_h_smooth_PP_a030.dat};
      \addplot[pline=4] table[x=h, y=eE4]{\data/cyl_h_smooth_PP_a030.dat};
      \addplot[pline=5] table[x=h, y=eE5]{\data/cyl_h_smooth_PP_a030.dat};
      \addplot[pline=6] table[x=h, y=eE6]{\data/cyl_h_smooth_PP_a030.dat};
      \addplot[pline=7] table[x=h, y=eE7]{\data/cyl_h_smooth_PP_a030.dat};
      \addplot[pline=8] table[x=h, y=eE8]{\data/cyl_h_smooth_PP_a030.dat};
      \nextgroupplot[title={\PM, Lam\'e}, ylabel={$\eE$}]
      \addplot[pline=2] table[x=h, y=eE2]{\data/cyl_h_lame_PM_a030.dat};
      \addplot[pline=3] table[x=h, y=eE3]{\data/cyl_h_lame_PM_a030.dat};
      \addplot[pline=4] table[x=h, y=eE4]{\data/cyl_h_lame_PM_a030.dat};
      \addplot[pline=5] table[x=h, y=eE5]{\data/cyl_h_lame_PM_a030.dat};
      \addplot[pline=6] table[x=h, y=eE6]{\data/cyl_h_lame_PM_a030.dat};
      \addplot[pline=7] table[x=h, y=eE7]{\data/cyl_h_lame_PM_a030.dat};
      \addplot[pline=8] table[x=h, y=eE8]{\data/cyl_h_lame_PM_a030.dat};
      \nextgroupplot[title={\PH, Lam\'e}]
      \addplot[pline=2] table[x=h, y=eE2]{\data/cyl_h_lame_PH_a030.dat};
      \addplot[pline=3] table[x=h, y=eE3]{\data/cyl_h_lame_PH_a030.dat};
      \addplot[pline=4] table[x=h, y=eE4]{\data/cyl_h_lame_PH_a030.dat};
      \addplot[pline=5] table[x=h, y=eE5]{\data/cyl_h_lame_PH_a030.dat};
      \addplot[pline=6] table[x=h, y=eE6]{\data/cyl_h_lame_PH_a030.dat};
      \addplot[pline=7] table[x=h, y=eE7]{\data/cyl_h_lame_PH_a030.dat};
      \addplot[pline=8] table[x=h, y=eE8]{\data/cyl_h_lame_PH_a030.dat};
      \nextgroupplot[title={\PP, Lam\'e}]
      \addplot[pline=2] table[x=h, y=eE2]{\data/cyl_h_lame_PP_a030.dat};
      \addplot[pline=3] table[x=h, y=eE3]{\data/cyl_h_lame_PP_a030.dat};
      \addplot[pline=4] table[x=h, y=eE4]{\data/cyl_h_lame_PP_a030.dat};
      \addplot[pline=5] table[x=h, y=eE5]{\data/cyl_h_lame_PP_a030.dat};
      \addplot[pline=6] table[x=h, y=eE6]{\data/cyl_h_lame_PP_a030.dat};
      \addplot[pline=7] table[x=h, y=eE7]{\data/cyl_h_lame_PP_a030.dat};
      \addplot[pline=8] table[x=h, y=eE8]{\data/cyl_h_lame_PP_a030.dat};
    \end{groupplot}
  \end{tikzpicture}\\[1mm]
  \ref{leg:p}
  \caption{Thick cylinder, $\alpha=30^\circ$ (no folded elements): energy error versus element size
    for the analytic field (top) and the Lam\'e field (bottom).}
  \label{fig:cyl-h}
\end{figure}

\begin{table}[t]
  \centering\small
  \caption{Thick cylinder: observed rates of $\eE$ between $N=4$ and $N=8$ (optimal: $p$);
    ``--'': error at round-off level.}
  \label{tab:rates}
  \begin{tabular}{@{}llrrrrrrr@{}}
\toprule
 & & $p=2$ & $p=3$ & $p=4$ & $p=5$ & $p=6$ & $p=7$ & $p=8$\\
\midrule
analytic, $\alpha=15^\circ$ & \PM & 2.00 & 2.99 & 3.99 & 4.99 & 5.98 & -- & --\\
 & \PH & 2.00 & 3.00 & 4.00 & 4.99 & 6.01 & -- & --\\
 & \PP & 1.08 & 1.23 & 2.21 & 2.03 & 2.94 & 3.39 & 4.05\\
\addlinespace[2pt]
analytic, $\alpha=30^\circ$ & \PM & 2.00 & 2.86 & -0.98 & 2.40 & -- & -- & --\\
 & \PH & 2.00 & 3.00 & 4.01 & 5.00 & 6.02 & -- & --\\
 & \PP & 1.02 & 1.06 & 2.01 & 1.98 & 2.98 & 3.11 & 3.94\\
\addlinespace[2pt]
Lam\'e, $\alpha=15^\circ$ & \PM & 1.97 & 2.91 & 3.90 & 4.81 & 5.83 & -- & --\\
 & \PH & 1.96 & 2.94 & 3.94 & 4.92 & 5.91 & 6.88 & 7.87\\
 & \PP & 1.07 & 1.43 & 2.16 & 2.81 & 3.49 & 4.19 & 4.78\\
\addlinespace[2pt]
Lam\'e, $\alpha=30^\circ$ & \PM & 1.97 & 2.94 & 0.76 & 2.99 & 4.56 & -- & --\\
 & \PH & 1.96 & 2.96 & 3.94 & 4.94 & 5.92 & 6.90 & 7.87\\
 & \PP & 1.03 & 1.21 & 2.07 & 2.51 & 3.20 & 3.85 & 4.42\\
\bottomrule
\end{tabular}

\end{table}

\begin{figure}[t]
  \centering
  \begin{tikzpicture}
    \begin{groupplot}[group style={group size=3 by 1, horizontal sep=3mm, yticklabels at=edge left},
        width=5.0cm, height=4.6cm, ymode=log, xmin=2, xmax=8, ymin=1e-14, ymax=1e-1,
        xlabel={$p$}, unbounded coords=jump]
      \nextgroupplot[title={$\alpha=15^\circ$}, ylabel={$e_n$ ($N=8$)},
        legend style={font=\scriptsize, at={(0.03,0.03)}, anchor=south west}]
      \addplot[sPP] table[x=p, y=eNN8]{\data/cyl_p_smooth_PP_a015.dat};
      \addplot[sPM] table[x=p, y=eNN8]{\data/cyl_p_smooth_PM_a015.dat};
      \addplot[sPH] table[x=p, y=eNN8]{\data/cyl_p_smooth_PH_a015.dat};
      \legend{PP, PM, PH}
      \nextgroupplot[title={$\alpha=30^\circ$}]
      \addplot[sPP] table[x=p, y=eNN8]{\data/cyl_p_smooth_PP_a030.dat};
      \addplot[sPM] table[x=p, y=eNN8]{\data/cyl_p_smooth_PM_a030.dat};
      \addplot[sPH] table[x=p, y=eNN8]{\data/cyl_p_smooth_PH_a030.dat};
      \nextgroupplot[title={$\alpha=45^\circ$ (folded)}]
      \addplot[sPP] table[x=p, y=eNN8]{\data/cyl_p_smooth_PP_a045.dat};
      \addplot[sPM] table[x=p, y=eNN8]{\data/cyl_p_smooth_PM_a045.dat};
      \addplot[sPH] table[x=p, y=eNN8]{\data/cyl_p_smooth_PH_a045.dat};
    \end{groupplot}
  \end{tikzpicture}
  \caption{Thick cylinder, analytic field: $p$-convergence of the nodal error on the finest
    mesh.}
  \label{fig:cyl-p}
\end{figure}

\paragraph{Approximation versus stability.} Inserting the exact nodal values into the trial
functions gives the interpolation error of each space; the ratio of discretisation to
interpolation error isolates the stability of the discrete problem (\cref{tab:quasiopt}). For \PH\ it
lies between $0.78$ and $1.00$ in all $158$ cases on the meshes without folds with errors above the
round-off level, for \PP\ between $0.84$ and $1.00$ ($165$ cases): both are quasi-optimal with a
constant close to one. For \PM\ it ranges from $0.07$ to $10^5$; at $\alpha=30^\circ$, $p\ge3$ the
discrete \PM\ solution is up to $64$ times worse \emph{or} up to five times better than its
interpolant, i.e.\ it does not inherit the approximation property of its space. The \PM\ space
itself also approximates worse where its Lebesgue constant is large: at $\alpha=30^\circ$, $p\le6$, the
\PM\ interpolant of the analytic field is up to nine times less accurate than the \PH\
interpolant.

\begin{table}[t]
  \centering\small
  \caption{Thick cylinder, $N=8$: energy error of the discrete solution divided by the energy
    error of the interpolant of the exact field in the same trial space.}
  \label{tab:quasiopt}
  \begin{tabular}{@{}llrrrrrrr@{}}
\toprule
 & & $p=2$ & $p=3$ & $p=4$ & $p=5$ & $p=6$ & $p=7$ & $p=8$\\
\midrule
analytic, $\alpha=15^\circ$ & \PM & 0.99 & 0.98 & 1.00 & 0.98 & 1.01 & -- & --\\
 & \PH & 0.99 & 0.97 & 0.98 & 0.96 & 0.98 & -- & --\\
 & \PP & 0.97 & 0.94 & 0.96 & 0.93 & 0.96 & 0.95 & 0.96\\
\addlinespace[2pt]
analytic, $\alpha=30^\circ$ & \PM & 0.91 & 0.20 & 64.38 & 3.40 & 0.38 & -- & --\\
 & \PH & 0.96 & 0.93 & 0.93 & 0.89 & 0.90 & -- & --\\
 & \PP & 0.94 & 0.92 & 0.92 & 0.89 & 0.93 & 0.93 & 0.93\\
\addlinespace[2pt]
Lam\'e, $\alpha=15^\circ$ & \PM & 1.00 & 0.99 & 0.99 & 0.98 & 0.99 & -- & --\\
 & \PH & 0.98 & 0.95 & 0.96 & 0.95 & 0.96 & 0.94 & 0.95\\
 & \PP & 0.98 & 0.96 & 0.98 & 0.97 & 0.98 & 0.97 & 0.98\\
\addlinespace[2pt]
Lam\'e, $\alpha=30^\circ$ & \PM & 0.99 & 0.48 & 14.98 & 5.46 & 8.75 & -- & --\\
 & \PH & 0.93 & 0.91 & 0.92 & 0.89 & 0.90 & 0.88 & 0.89\\
 & \PP & 0.95 & 0.93 & 0.96 & 0.93 & 0.95 & 0.94 & 0.95\\
\bottomrule
\end{tabular}

\end{table}

\paragraph{Lebesgue constants, conditioning and spectra.} \Cref{tab:lebesgue} lists the
maxima over the elements of the finest mesh. The Lebesgue constant of \PM\ reaches $10^3$ at
$30^\circ$, $p=4$ and $1.3\cdot10^8$ at $45^\circ$, $p=8$; that of \PH\ stays within $30\,\%$ of the GLL value
up to $15^\circ$ and $\le95$ at $30^\circ$. $\kappa_2(\bm V)$ of \PM\ exceeds its Lebesgue constant by up
to two orders of magnitude (extrapolation at the reference points, \cref{prop:kappa}); that of \PH\
is $10^2$--$10^5$ times smaller than that of \PM\ at $30^\circ$, $p\ge6$. At mesh level, the smallest singular values of the global
stiffness matrices of the three elements agree to $1\,\%$ (14 cases, $\alpha\le30^\circ$), but the
largest singular value of \PM\ exceeds that of \PP\ by up to a factor $344$ ($30^\circ$, $N=4$,
$p=6$), whereas that of \PH\ stays within $1.2$ of it. In the expansion of the error of the analytic
field in the singular triplets of $\bm K$, a single one of the six lowest triplets carries
$21$--$68\,\%$ of the \PM\ error in four of five cases with $\alpha\ge30^\circ$; for \PH\ and \PP\ the error
is spread over the spectrum (largest single share $2$--$8\,\%$), and the lowest $50$ singular values
of \PH\ agree with those of \PP\ to $2.3\,\%$.

\begin{table}[t]
  \centering\small
  \caption{Thick cylinder, $N=8$: Lebesgue constant $\Lambda$ ($16\times16$ sample points per
    element, a lower bound) and condition number $\kappa_2(\bm V)$ of the interpolation matrix,
    maximum over the elements.}
  \label{tab:lebesgue}
  \begin{tabular}{@{}llrrrrrrr@{}}
\toprule
 & & $p=2$ & $p=3$ & $p=4$ & $p=5$ & $p=6$ & $p=7$ & $p=8$\\
\midrule
$\Lambda$, $\alpha=15^\circ$ & \PM & $2.5$ & $6.0$ & $1.5\cdot10^{1}$ & $3.7\cdot10^{1}$ & $1.3\cdot10^{2}$ & $1.2\cdot10^{4}$ & $1.3\cdot10^{4}$\\
 & \PH & $1.8$ & $2.4$ & $2.8$ & $3.1$ & $3.5$ & $4.5$ & $4.2$\\
$\Lambda$, $\alpha=30^\circ$ & \PM & $1.0\cdot10^{1}$ & $2.4\cdot10^{2}$ & $1.1\cdot10^{3}$ & $1.6\cdot10^{3}$ & $6.2\cdot10^{3}$ & $4.2\cdot10^{6}$ & $1.2\cdot10^{6}$\\
 & \PH & $2.0$ & $4.8$ & $7.5$ & $1.6\cdot10^{1}$ & $2.5\cdot10^{1}$ & $7.4\cdot10^{1}$ & $9.5\cdot10^{1}$\\
$\Lambda$, $\alpha=45^\circ$ & \PM & $1.9\cdot10^{2}$ & $1.5\cdot10^{3}$ & $1.3\cdot10^{4}$ & $5.0\cdot10^{5}$ & $8.3\cdot10^{5}$ & $1.4\cdot10^{6}$ & $1.3\cdot10^{8}$\\
 & \PH & $2.8$ & $1.0\cdot10^{1}$ & $2.1\cdot10^{1}$ & $4.7\cdot10^{1}$ & $9.1\cdot10^{1}$ & $2.6\cdot10^{2}$ & $5.5\cdot10^{2}$\\
 & GLL & $1.56$ & $2.22$ & $2.68$ & $3.00$ & $3.41$ & $3.47$ & $4.15$\\
\midrule
$\kappa_2(\bm V)$, $\alpha=15^\circ$ & \PM & $4.8$ & $1.2\cdot10^{1}$ & $6.0\cdot10^{1}$ & $2.9\cdot10^{2}$ & $1.9\cdot10^{3}$ & $3.1\cdot10^{5}$ & $2.4\cdot10^{5}$\\
 & \PH & $5.5$ & $9.7$ & $1.9\cdot10^{1}$ & $4.0\cdot10^{1}$ & $9.4\cdot10^{1}$ & $2.3\cdot10^{2}$ & $6.0\cdot10^{2}$\\
$\kappa_2(\bm V)$, $\alpha=30^\circ$ & \PM & $3.8\cdot10^{1}$ & $1.7\cdot10^{3}$ & $5.9\cdot10^{3}$ & $8.5\cdot10^{3}$ & $8.6\cdot10^{4}$ & $2.6\cdot10^{8}$ & $2.6\cdot10^{8}$\\
 & \PH & $9.0$ & $2.4\cdot10^{1}$ & $6.0\cdot10^{1}$ & $1.8\cdot10^{2}$ & $7.0\cdot10^{2}$ & $2.7\cdot10^{3}$ & $9.8\cdot10^{3}$\\
$\kappa_2(\bm V)$, $\alpha=45^\circ$ & \PM & $8.4\cdot10^{2}$ & $1.5\cdot10^{4}$ & $2.1\cdot10^{5}$ & $2.9\cdot10^{7}$ & $2.8\cdot10^{8}$ & $4.7\cdot10^{9}$ & $7.9\cdot10^{11}$\\
 & \PH & $1.8\cdot10^{1}$ & $9.1\cdot10^{1}$ & $4.1\cdot10^{2}$ & $2.4\cdot10^{3}$ & $3.3\cdot10^{4}$ & $3.2\cdot10^{5}$ & $4.0\cdot10^{6}$\\
\bottomrule
\end{tabular}

\end{table}

\paragraph{Limits of the hybrid space.} (i) On the folded meshes with $\alpha\ge60^\circ$ the
interpolation matrix of \PH\ becomes singular in $6$ of the $24$ cases with $p\ge6$ (that of \PM\ in
$7$). \Cref{tab:unisolvence} confirms \cref{prop:factor}: $\mathrm{rcond}(\bm A)$ follows
$\mathrm{rcond}(\bm V_T)$ within a factor $30$, the determinant identity~\eqref{eq:factor} holds to
$\le3\cdot10^{-13}$ (relative) for $\mathrm{rcond}(\bm V_T)\ge5\cdot10^{-6}$ and degrades with the
conditioning as expected for computed determinants, and the singular cases are exactly those with
$\mathrm{rcond}(\bm A)\le4\cdot10^{-16}$ (all other cases: $\ge2\cdot10^{-15}$). (ii) On the same meshes the \PH\ error for $p=2$ ($60^\circ$)
and $p\le4$ ($90^\circ$) is non-monotone in $N$ and of order $10^{-2}$--$1$, as for the other two
elements; the singular complement integrand (\cref{sec:hybrid}) is one possible cause. (iii) The
error constant of \PH\ depends on the distortion: for the Lam\'e field at $\alpha=15^\circ$, where \PM\ is
stable, \PH\ is up to $4.1$ times less accurate than \PM\ ($p\le7$), because its interpolation error
is $1.2$--$5.7$ times larger ($N=8$); the parametric complement approximates the $r^{-2}$ terms of this
field worse than the full $Q_p$ in skew coordinates.

\begin{table}[t]
  \centering\small
  \caption{Unisolvence of the hybrid space on the thick-cylinder meshes: minimum over the elements
    of $\mathrm{rcond}$ of $\bm V_T$, of $\bm A$ (\cref{prop:factor}) and of $\bm V$ of
    $Q_p(\hat x,\hat y)$; number of elements with $\mathrm{rcond}(\bm V_T)<\varepsilon$.}
  \label{tab:unisolvence}
  \begin{tabular}{@{}rrrcccc@{}}
\toprule
$\alpha$ & $N$ & $p$ & $\min\mathrm{rcond}(\bm V_T)$ & $\min\mathrm{rcond}(\bm A)$ & $\min\mathrm{rcond}(\bm V_Q)$ & singular\\
\midrule
$0^\circ$ & 8 & 8 & $1.9\cdot10^{-3}$ & $5.5\cdot10^{-2}$ & $1.0\cdot10^{-2}$ & 0\\
$15^\circ$ & 8 & 8 & $7.5\cdot10^{-5}$ & $5.0\cdot10^{-4}$ & $4.0\cdot10^{-7}$ & 0\\
$30^\circ$ & 8 & 6 & $1.0\cdot10^{-4}$ & $5.5\cdot10^{-4}$ & $2.9\cdot10^{-6}$ & 0\\
$30^\circ$ & 8 & 7 & $2.1\cdot10^{-5}$ & $1.1\cdot10^{-4}$ & $6.3\cdot10^{-10}$ & 0\\
$30^\circ$ & 8 & 8 & $5.3\cdot10^{-6}$ & $2.7\cdot10^{-5}$ & $6.3\cdot10^{-10}$ & 0\\
$45^\circ$ & 8 & 6 & $2.8\cdot10^{-6}$ & $1.3\cdot10^{-5}$ & $1.2\cdot10^{-9}$ & 0\\
$45^\circ$ & 8 & 8 & $1.8\cdot10^{-8}$ & $8.1\cdot10^{-8}$ & $4.7\cdot10^{-13}$ & 0\\
$60^\circ$ & 4 & 8 & $1.4\cdot10^{-17}$ & $9.8\cdot10^{-17}$ & $7.5\cdot10^{-19}$ & 4\\
$60^\circ$ & 8 & 6 & $3.1\cdot10^{-16}$ & $2.2\cdot10^{-15}$ & $1.3\cdot10^{-16}$ & 0\\
$60^\circ$ & 8 & 7 & $5.2\cdot10^{-19}$ & $6.5\cdot10^{-18}$ & $1.7\cdot10^{-18}$ & 8\\
$60^\circ$ & 8 & 8 & $3.4\cdot10^{-20}$ & $3.1\cdot10^{-19}$ & $6.1\cdot10^{-21}$ & 8\\
$90^\circ$ & 1 & 7 & $7.8\cdot10^{-16}$ & $7.8\cdot10^{-15}$ & $1.5\cdot10^{-15}$ & 0\\
$90^\circ$ & 1 & 8 & $6.1\cdot10^{-18}$ & $6.1\cdot10^{-17}$ & $5.3\cdot10^{-19}$ & 1\\
$90^\circ$ & 8 & 6 & $1.2\cdot10^{-14}$ & $9.3\cdot10^{-14}$ & $6.0\cdot10^{-15}$ & 0\\
$90^\circ$ & 8 & 7 & $3.9\cdot10^{-17}$ & $3.9\cdot10^{-16}$ & $2.8\cdot10^{-17}$ & 8\\
$90^\circ$ & 8 & 8 & $3.9\cdot10^{-20}$ & $9.4\cdot10^{-19}$ & $1.2\cdot10^{-19}$ & 16\\
\bottomrule
\end{tabular}

\end{table}

\section{Discussion and conclusions}\label{sec:conclusions}
The metric trial functions of unsymmetric spectral elements suffer from two distinct
deficiencies that grow with the polynomial degree and the element distortion: an
ill-conditioned interpolation \emph{system} and an ill-conditioned interpolation
\emph{problem}. The first is a matter of the basis, the second of the trial space.

\paragraph{Basis.} The trial space, its Lebesgue constant, distortion immunity and the exact
element matrix are basis invariant (\cref{prop:invariance}). The scaled tensor-product Lagrange
basis on GLL points makes the interpolation matrix a permutation on affine elements, so the
unsymmetric element coincides with the conventional one to round-off there (symmetric
$\bm K^e$ for all $p$), and on distorted elements it reduces $\kappa(\bm V)$ to the Lebesgue
function of the interpolation at the reference points (\cref{prop:kappa}). Compared with
monomials it lowers $\kappa$ by $4.6\cdot10^3$--$2.9\cdot10^4$ and the error of the computed shape
functions by $3.6\cdot10^2$--$8.2\cdot10^3$ on the distorted test elements at $p=16$, and it removes the
spurious rejection of elements by a singularity check (five element--degree combinations on the panel meshes). Its
benefit is confined to round-off: approximation, the amplification of nodal perturbations and
the stability of the discrete problem are unchanged, and the static results of the thick cylinder
change only at round-off level. A better basis \emph{cannot} improve the Lebesgue constant.

\paragraph{Hybrid space.} Completeness requires metric functions only for $P_p(\bx)$. The hybrid
space $T_p$ (\cref{def:hybrid}) keeps them for $P_p$ and represents the remaining
$\tfrac12p(p+1)$ functions in parametric coordinates. It retains distortion immunity, frame
invariance and interpolation, reduces to the conventional spectral element on straight-sided
elements, and its interpolant is the conventional one plus the interpolation error of a
polynomial composed with the element map (\cref{prop:structure}), which bounds its Lebesgue
constant by the GLL value plus a term that vanishes for bilinear maps. Unisolvence is equivalent
to the non-singularity of a $d\times d$ matrix $\bm A$ (\cref{prop:factor}). Numerically, the hybrid
space
\begin{itemize}
\item keeps the Lebesgue constant at or near the GLL value on single elements up to $p=16$
  ($\le8.5$, against $5.5\cdot10^4$--$4.0\cdot10^7$ for $Q_p$ in skew coordinates) and $\le95$ on the
  thick-cylinder meshes without folds (against up to $4.2\cdot10^6$);
\item reproduces polynomial fields to $\le10^{-9}$ on all meshes, including folded ones, where
  $Q_p$ in skew coordinates is limited to $10^{-5}$--$10^{-2}$ by round-off;
\item produces no eigenvalues with negative real part and imaginary parts $\le1.8\,\%$ of the
  modulus, where $Q_p$ in skew coordinates produces up to $25$ unstable modes on curved meshes;
\item converges with the optimal order $p$ and quasi-optimally (discrete error $0.78$--$1.00$ times
  the interpolation error) on all non-folded meshes of a mesh sequence with fixed relative
  distortion, on which $Q_p$ in skew coordinates is optimal only for moderate distortion and the
  conventional element converges with $\approx\lfloor p/2\rfloor$; the gain over the conventional
  element reaches $2.6\cdot10^4$ in the energy norm ($p=7$, $\alpha=30^\circ$, $N=8$).
\end{itemize}

\paragraph{Limitations.} (i) No stability theory (discrete inf--sup
condition~\cite{Babuska1971Error}) is available for either unsymmetric element; quasi-optimality
of \PH\ is observed, not proven. (ii) Unisolvence of $T_p$ is guaranteed only near affine maps; it
fails on some folded elements at $p\ge7$, and the complement part of the element matrix has a
singular integrand where $\det\bm J$ changes sign; on folded meshes \PH\ is reliable only for
$p\ge4$--$5$. (iii) The error constant of \PH\ depends on the distortion: where $Q_p$ in skew
coordinates is stable, it can approximate better (by factors up to $5.7$ for the Lam\'e field at
$\alpha=15^\circ$, up to $17$ for $e^x\cos y$ on curved meshes). (iv) The studies are two-dimensional,
linear-elastic and restricted to $p\le16$ (element level) and $p\le8$ (convergence); 3D elements and
non-linear problems remain to be examined.

For high-order unsymmetric spectral elements we therefore recommend the scaled Lagrange basis
for any metric space, and the hybrid space $T_p$ whenever $p\gtrsim4$, the distortion is not
small, or the spectrum matters (transient and modal analysis).

\section*{Reproduction}\label{sec:reproduction}
Code base: USEM-MATLAB after commit \texttt{e293da1} (\texttt{ShapeFunctions2D\_Metric},
\texttt{p.metricBasis}, \texttt{p.metricMethod}). Element and mesh level:
\texttt{studies/metric\_basis/*.m}; patch tests: \texttt{studies/rajendran\_panel};
thick cylinder: \texttt{studies/thick\_cylinder}; unisolvence:
\texttt{studies/manuscript\_metric/matlab/unisolvence\_check.m}; tables and data of this paper:
\texttt{studies/manuscript\_metric/make\_data.py}.


\bibliographystyle{plainnat}
\bibliography{../../literature/references,refs}
\end{document}
